\documentclass[a4paper, 12pt]{article}
\input{../preamble.tex}


\begin{document}

\begin{center}
{\large КОНСПЕКТ ЛЕКЦИЙ}

ПО ДИСЦИПЛИНЕ <<АНАЛИТИЧЕСКАЯ ГЕОМЕТРИЯ>>

\tableofcontents
\label{contents}

\pagebreak

\end{center}
\vspace{10pt}

\section{Матрицы}

Матрицы возникают как средство решения линейных систем.

\textbf{Пример.} Предприятие использует $m$ ресурсов в объёмах
$b_1, \dots, b_m$, выпускает $n$ продуктов в объёме $x_1, \dots, x_n$.
$a_{ij}$ --- расход ресурса $i$ на единицу продукта $j$, $i = \overline{1, m}$
(от 1 до $m$). Тогда $x_j$ удовлетворяют системе:
\[
  \begin{cases}
    a_{11} x_1 + a_{12} x_2 + \dots + a_{1n} x_n = b_1,\\
    \dotfill\\
    a_{m1} x_1 + a_{m2} x_2 + \dots + a_{mn} x_n = b_m.
  \end{cases}
\]
Пояснить: $a_{m2} x_2$ --- использовано ресурса $m$ на выпуск продукта №\,2
объёма $x_2$.

\textbf{Идея:} решить систему, используя таблицу её коэффициентов.

\textbf{Определение.} Матрица $A$ размерности $n \times m$ --- это таблица,
в которой $n$ строк, $m$ столбцов, с элементами $a_{ij}$
($i$ --- номер строки, $j$ --- номер столбца).

Пр.:
\[
  A = \begin{pmatrix} 2 & 3 & 0 \\ -1 & 4 & 5 \end{pmatrix} \quad (2 \times 3),
  \qquad a_{13} = 0, \quad a_{22} = 4.
\]

\subsection{Частные случаи}

\begin{enumerate}[label=\arabic*)]
  \item \textbf{Строка (вектор-строка):} $a = \begin{pmatrix} 1 & 2 & 3 \end{pmatrix}$
    --- $1 \times 3$. В общем случае $1 \times n$.
  \item \textbf{Столбец (вектор-столбец).} Далее вектор --- всегда столбец,
    то есть упорядоченный набор чисел, или матрица $m \times 1$:
    \[
      a = \begin{pmatrix} -1 \\ 2 \\ 0 \end{pmatrix} \quad (3 \times 1).
    \]
    Матрицы будем обозначать большими буквами, векторы --- малыми буквами.
  \item \textbf{Треугольная матрица} --- все элементы выше (ниже) главной
    диагонали равны 0: $a_{ij} = 0$, если $i < j$ ($i > j$).

    Пр.: верхняя треугольная матрица
    \[
      a = \begin{pmatrix} 1 & 4 & 5 \\ 0 & 2 & 6 \\ 0 & 0 & 3 \end{pmatrix};
    \]
    диагональ --- элементы вида $a_{ii}$; элементы под диагональю --- нули.
  \item \textbf{Диагональная матрица.} $A$ диагональная, если $a_{ij} = 0$
    при $i \ne j$:
    \[
      \begin{pmatrix} 0 & 0 & 0 & 0 \\ 0 & 1 & 0 & 0 \\ 0 & 0 & 2 & 0 \end{pmatrix}.
    \]
    Диагональные элементы могут быть нулями. % в оригинале: «# i≠j могут быть 0»
\end{enumerate}

\textbf{!} \underline{Единичная матрица} $E = \begin{pmatrix} 1 & & 0 \\ & \ddots & \\ 0 & & 1 \end{pmatrix}$
--- $n \times n$. Это и диагональная, и треугольная матрица.

\textbf{!} \underline{Нулевая матрица} $O_{nm} = \begin{pmatrix} 0 & \dots & 0 \\ \vdots & & \vdots \\ 0 & \dots & 0 \end{pmatrix}$
(все нули) --- и треугольная, и диагональная.

\subsection{Операции над матрицами}

\begin{enumerate}[label=\arabic*)]
  \item \textbf{Транспонирование.} При этом каждая строка №\,$i$ матрицы $A$
    становится столбцом №\,$i$ матрицы $A^T$ (строки и столбцы <<меняются
    местами>>). При этом $A$ <<поворачивается>> относительно диагонали.

    Пр.:
    \[
      A = \begin{pmatrix} 1 & 2 & 3 & 4 \\ 5 & 6 & 7 & 8 \end{pmatrix}^T
        = \begin{pmatrix} 1 & 5 \\ 2 & 6 \\ 3 & 7 \\ 4 & 8 \end{pmatrix}
    \]
    (диагональные элементы остаются на месте).

  \item \textbf{Линейные операции.}
    \begin{enumerate}[label=\alph*)]
      \item Умножение на число: $\alpha \cdot A = (\alpha a_{ij})$. Т.\,е.
        при умножении матрицы на число каждый её элемент умножается на число.
      \item Сложение матриц: $A + B = C$, где $A = (a_{ij})$, $B = (b_{ij})$,
        $C = (c_{ij})$, $c_{ij} = a_{ij} + b_{ij}$. При сложении матриц
        складываем их соответствующие элементы $\Rightarrow$ матрицы должны
        быть одинаковой размерности.
    \end{enumerate}

  \item \textbf{Скалярное произведение векторов} (алгебраическое обозначение).
    \[
      a = \begin{pmatrix} a_1 \\ a_2 \\ \vdots \\ a_n \end{pmatrix}, \qquad
      b = \begin{pmatrix} b_1 \\ b_2 \\ \vdots \\ b_n \end{pmatrix},
    \]
    \begin{align*}
      a^T b &= \begin{pmatrix} a_1 & a_2 & \dots & a_n \end{pmatrix} \cdot
        \begin{pmatrix} b_1 \\ b_2 \\ \vdots \\ b_n \end{pmatrix} \\
      &= \sum_{i=1}^{n} a_i b_i = a_1 b_1 + a_2 b_2 + \dots + a_n b_n
    \end{align*}
    --- сумма произведений соответствующих координат.

    \textbf{Замечание.} С геометрической точки зрения
    $a^T b = 0 \Longleftrightarrow$ <<$a \perp b$>>.
    $a^2 = a^T a$ (скалярный квадрат); $|a| = \sqrt{a^2}$ (евклидова длина $a$).

  \item \textbf{Произведение матриц.}
    \[
      C = A \cdot B, \qquad
      \underbrace{(c_{ij})}_{n \times m} = \underbrace{(a_{ij})}_{n \times k}
      \cdot \underbrace{(b_{ij})}_{k \times m},
    \]
    если каждый элемент $c_{ij}$ --- это скалярное произведение строки №\,$i$
    матрицы $A$ на столбец №\,$j$ матрицы $B$.

    $\Downarrow$ Столбцов в матрице слева должно быть столько же, сколько
    строк в матрице справа.

    Пр.:
    \begin{enumerate}[label=\arabic*)]
      \item $\begin{pmatrix} 1 \\ 2 \\ 3 \end{pmatrix} \cdot \begin{pmatrix} 4 & 5 & 6 \end{pmatrix}
        = \begin{pmatrix} 1\cdot4 & 1\cdot5 & 1\cdot6 \\ 2\cdot4 & 2\cdot5 & 2\cdot6 \\ 3\cdot4 & 3\cdot5 & 3\cdot6 \end{pmatrix}$
        --- $3 \times 1$ на $1 \times 3$: можно перемножать, должна
        получиться матрица $3 \times 3$.
      \item $\begin{pmatrix} 1 \\ 2 \end{pmatrix} \cdot \begin{pmatrix} 3 \\ 4 \end{pmatrix}$
        --- $2 \times 1$ на $2 \times 1$: размерности не согласованы,
        перемножить нельзя.
      \item $\begin{pmatrix} 1 & 2 \\ 3 & 4 \end{pmatrix} \cdot
        \begin{pmatrix} 5 & 7 & 9 \\ 6 & 8 & 10 \end{pmatrix}$
        --- $2 \times 2$ на $2 \times 3$: можно умножать по принципу
        <<строка на столбец>>:
        \[
          = \begin{pmatrix}
              1\cdot5 + 2\cdot6 & 1\cdot7 + 2\cdot8 & 1\cdot9 + 2\cdot10 \\
              3\cdot5 + 4\cdot6 & 3\cdot7 + 4\cdot8 & 3\cdot9 + 4\cdot10
            \end{pmatrix} \quad (2 \times 3).
        \]
    \end{enumerate}

    \textbf{!} $B \cdot A$ --- нельзя ($2 \times 3$ на $2 \times 2$, $3 \ne 2$).
    В общем случае $AB \ne BA$. Если $AB = BA$, то $A$ и $B$ коммутируют.
\end{enumerate}

\subsection{Свойства операций}

В основном аналогичны операциям с числами. Эти свойства непосредственно
следуют из определений.

\textbf{Свойства сложения матриц:}
\begin{itemize}[noitemsep]
  \item $A + B = B + A$ --- коммутативность;
  \item $A + (-A) = O_{nm}$ --- наличие противоположного элемента;
  \item $A + O_{nm} = A$ --- нейтральность нулевой матрицы при сложении;
  \item $A + (B + C) = (A + B) + C$ --- ассоциативность;
  \item $(A + B)^T = A^T + B^T$ --- транспонирование суммы.
\end{itemize}

\textbf{Свойства умножения матрицы на число:}
\begin{itemize}[noitemsep]
  \item $\alpha (A + B) = \alpha A + \alpha B$ --- дистрибутивность;
  \item $(\alpha + \beta) A = \alpha A + \beta A$ --- дистрибутивность;
  \item $(\alpha \beta) A = \alpha (\beta A)$ --- ассоциативность.
\end{itemize}

\textbf{Свойства умножения матриц:}
\begin{itemize}[noitemsep]
  \item $C (A + B) = CA + CB$, \ $(A + B) C = AC + BC$ --- дистрибутивность;
    % В оригинале: C(A+B) = (A+B)C = CA + CB — неверно, C(A+B) и (A+B)C в общем случае различны.
  \item $A (BC) = (AB) C$ --- ассоциативность;
  \item $\alpha (AB) = (\alpha A) B = A (\alpha B)$ --- ассоциативность;
  \item $AE = EA = A$ --- нейтральность единичной матрицы при умножении;
  \item $(AB)^T = B^T A^T$ --- транспонирование произведения.
\end{itemize}

\subsection{Определитель матрицы (детерминант)}

\begin{enumerate}[label=\arabic*)]
  \item $A = (a_{11})$ --- $1 \times 1$ $\Rightarrow$ $\det A = a_{11}$.
  \item $A = \begin{pmatrix} a_{11} & a_{12} \\ a_{21} & a_{22} \end{pmatrix}$ --- $2 \times 2$
    $\Rightarrow$
    \[
      \det A = a_{11} a_{22} - a_{12} a_{21}
    \]
    (произведение диагональных элементов минус произведение внедиагональных
    элементов).
\end{enumerate}

Определитель квадратной $n \times n$ матрицы --- скалярное произведение
$\forall$ строки (или столбца) на вектор алгебраических дополнений элементов
этой строки (столбца):
\begin{align*}
  \det A &= \sum_{j=1}^{n} a_{ij} \cdot A_{ij} && \text{(разложение по строке)},\\
  \det A &= \sum_{i=1}^{n} a_{ij} \cdot A_{ij} && \text{(разложение по столбцу)},
\end{align*}
где $A_{ij} = (-1)^{i+j} M_{ij}$ --- алгебраическое дополнение для $a_{ij}$,
$M_{ij}$ --- минор для $a_{ij}$: определитель после вычёркивания строки $i$
и столбца $j$, в которых был этот элемент. Последовательно берём элементы
строки (столбца).

\textbf{Пр.} $A$ --- $3 \times 3$; 1 способ: разложение по 3-му столбцу.
\begin{align*}
  \det \begin{pmatrix} 5 & -1 & 4 \\ 3 & 0 & 1 \\ 2 & 3 & -4 \end{pmatrix}
  &= 4 \cdot (-1)^{1+3} \begin{vmatrix} 3 & 0 \\ 2 & 3 \end{vmatrix}
   + 1 \cdot (-1)^{2+3} \begin{vmatrix} 5 & -1 \\ 2 & 3 \end{vmatrix} \\
  &\quad + (-4) \cdot (-1)^{3+3} \begin{vmatrix} 5 & -1 \\ 3 & 0 \end{vmatrix} \\
  &= 4 \cdot 9 - 1 \cdot 17 - 4 \cdot 3 = 36 - 17 - 12 = 7.
\end{align*}
% В тетради последнее слагаемое посчитано как −3 и ответ 16. Верно: (−4)·3 = −12, det = 7
% (проверка по правилу Саррюса: 0 − 2 + 36 − 0 − 12 − 15 = 7).

Для матриц 3-го порядка существует также:
\begin{enumerate}[label=\arabic*)]
  \item \textbf{Правило треугольников:}
    \[
      \det A = (a_{11} a_{22} a_{33} + a_{12} a_{23} a_{31} + a_{21} a_{32} a_{13})
             - (a_{13} a_{22} a_{31} + a_{12} a_{21} a_{33} + a_{23} a_{32} a_{11}).
    \]
  \item \textbf{Правило Саррюса} (справа дописываются первые два столбца):
    \[
      \det A =
      \left|\begin{array}{ccc|cc}
        a_{11} & a_{12} & a_{13} & a_{11} & a_{12} \\
        a_{21} & a_{22} & a_{23} & a_{21} & a_{22} \\
        a_{31} & a_{32} & a_{33} & a_{31} & a_{32}
      \end{array}\right.
    \]
    \[
      = a_{11} a_{22} a_{33} + a_{12} a_{23} a_{31} + a_{13} a_{21} a_{32}
      - a_{13} a_{22} a_{31} - a_{11} a_{23} a_{32} - a_{12} a_{21} a_{33}.
    \]
\end{enumerate}

\subsection{Свойства определителей}

\begin{enumerate}[label=\arabic*)]
  \item $|A^T| = |A|$ --- при транспонировании определитель сохраняется.
  \item $|AB| = |A| \cdot |B|$.
  \item Если все элементы строки (столбца) умножить на число, то определитель
    умножится на это число.

    \textbf{Следствия:}
    \begin{enumerate}[label=\alph*)]
      \item если матрица $A$ имеет нулевую строку (столбец), то $|A| = 0$;
      \item $|\lambda A| = \lambda^n |A|$.
    \end{enumerate}
  \item При перемене мест любых двух строк (столбцов) изменяется знак
    определителя.
  \item Для треугольной матрицы $\det A = \prod\limits_{i=1}^{n} a_{ii}$ ---
    это произведение диагональных элементов.
  \item Определитель не изменится, если совершить элементарное преобразование:
    к элементам какой-либо строки (столбца) прибавить любую другую строку
    (столбец), умноженную на число.

    \textbf{Следствие:} если матрица имеет две равные или пропорциональные
    строки (столбца), то её определитель равен~0.
\end{enumerate}

\subsection{Обоснование свойств определителей}

Если раскрыть определитель, то окажется, что
\begin{align}
  |A| &= \sum_{\sigma} (-1)^{|\sigma|}\, a_{1\sigma(1)}\, a_{2\sigma(2)} \dots a_{n\sigma(n)} \notag\\
      &= \sum_{\sigma} (-1)^{|\sigma|} \prod_{i=1}^{n} a_{i\sigma(i)} \tag{1}
\end{align}
--- сумма по всем перестановкам элементов множества
$\{1, 2, \dots, n\}$. Всего $n!$ вариантов и таково же число слагаемых.

$|\sigma|$ --- это число инверсий перестановки, т.\,е. количество случаев,
когда из $i < j$ следует $\sigma(i) > \sigma(j)$.

\textbf{Пример}: $\{1, 2, 3, 4\} \to \{4, 2, 1, 3\}$ (перестановка): инверсии
$(4,2), (4,1), (4,3), (2,1)$ $\Rightarrow$ перестановка <<чётная>>.
% В тетради выписаны две инверсии; на самом деле их четыре: (4,2), (4,1), (4,3), (2,1) — перестановка всё равно чётная.

\begin{enumerate}[label=\arabic*.]
  \item $A = \begin{pmatrix} a_{11} & a_{12} \\ a_{21} & a_{22} \end{pmatrix}$,
    $|A| = a_{11} a_{22} - a_{12} a_{21}$. По формуле (1):
    \[
      |A| = (-1)^0 a_{11} a_{22} + (-1)^1 a_{12} a_{21}
    \]
    (берём из строки 1 столбец 1 --- перестановка $\{1, 2\}$;
    $\{1, 2\} \to \{2, 1\}$ --- одна инверсия).
  \item Определитель порядка 3: по (1)
    \[
      |D| = \sum_{\sigma} (-1)^{|\sigma|}\, a_{1\sigma(1)} \dots \lambda a_{i\sigma(i)} \dots a_{n\sigma(n)}.
    \]
\end{enumerate}

Установим свойства определителей.
\begin{enumerate}[label=\arabic*)]
  \item Если строка (столбец) умножается на число $\lambda$, то оно выносится
    за знак определителя. Действительно, каждое слагаемое в (1) содержит
    ровно один множитель из этой строки $\Rightarrow$ $= \lambda |A|$.

    Если в определителе нулевая строка или столбец, то $\det = 0$.
    Действительно, можно считать, что вся строка имеет множитель
    $\lambda = 0$.
  \item $\det A^T = \det A$. Действительно, оба определителя вычисляются
    по формуле (1) и имеют одинаковые слагаемые:
    \[
      \det A^T = \sum_{\sigma} (-1)^{|\sigma|} \prod_{i=1}^{n} a_{\sigma(i) i}
    \]
    --- те же слагаемые, что и в (1). Самостоятельно проверить для матриц
    $2 \times 2$, $3 \times 3$.
  \item При перестановке двух строк (столбцов) $\det$ меняет знак.
    Действительно, рассмотрим перестановку двух строк: при этом все слагаемые
    останутся теми же, количество инверсий изменяется на 1 $\Rightarrow$
    все $(-1)^{|\sigma|}$ поменяют знак $\Rightarrow$ $-\det A$.

    \textbf{Следствие:} если в матрице две строки (столбца) пропорциональны
    или равны, то $\det = 0$. Действительно, без ограничения общности будем
    считать, что пропорциональны строки 1, 2:
    \[
      A = \begin{pmatrix}
        a_{11} & a_{12} & \dots & a_{1n} \\
        \lambda a_{11} & \lambda a_{12} & \dots & \lambda a_{1n} \\
        a_{31} & \dots & \dots & a_{3n} \\
        \vdots & & & \vdots \\
        a_{n1} & \dots & \dots & a_{nn}
      \end{pmatrix}.
    \]
    При вычислении определителя $\lambda$ выносим за знак определителя по
    свойству 1; полученный определитель обозначим $\det A'$:
    $\det A = \lambda \det A'$. Переставим строки 1, 2 $\Rightarrow$ должен
    поменяться знак, но по сути матрица без изменений, т.\,е.
    $\det A' = -\det A'$ $\Rightarrow$ $\det A' = 0$.
  \item \textbf{!} Если к строке (столбцу) прибавить строку (столбец),
    умноженную на число, то определитель не изменится. Действительно, без
    ограничения общности умножим строку 1 на $\lambda$, прибавим к строке 2
    и получим матрицу $B$. Разложим по строке 2:
    \begin{align*}
      \det B &= \sum_{j=1}^{n} (-1)^{2+j} (a_{2j} + \lambda a_{1j}) M_{2j} \\
      &= \underbrace{\sum_{j=1}^{n} (-1)^{2+j} a_{2j} M_{2j}}_{\det A}
      + \lambda \underbrace{\sum_{j=1}^{n} (-1)^{2+j} a_{1j} M_{2j}}_{= 0},
    \end{align*}
    т.\,к. второе слагаемое --- это определитель матрицы, в которой одинаковы
    строки 1, 2.
\end{enumerate}

В частном случае мы видим, что скалярное произведение строки на вектор
алгебраических дополнений любой другой строки равно~0. То же для столбцов.

\subsection{Обратная матрица}

\textbf{Идея:} для линейного уравнения $ax = b$ $\exists!$ решение, если
$a \ne 0$; тогда $x = \frac{1}{a} \cdot b$, где $\frac1a$ --- обратный элемент
для $a$ при умножении.

? Как обобщить эту идею на случай матричного уравнения
\[
  A \cdot X = B, \tag{1}
\]
где $A$, $B$ даны, $X$ --- найти.

Для ответа дадим определение: $A^{-1}$ --- \underline{обратная матрица}
для $A$, если
\[
  \boxed{A^{-1} A = A A^{-1} = E} \tag{2}
\]
($E$ --- единичная матрица, нейтральный элемент при умножении матриц).

\textbf{Замечание.} В математике любое определение должно быть
целесообразным --- предназначенным для решения какой-то задачи --- и
корректным, т.\,е. непротиворечивым и полным. % конец фразы в оригинале обрывается

? Целесообразно ли определение $A^{-1}$: $\exists A^{-1}$ $\Rightarrow$
умножим (1) на $A^{-1}$ слева $\Rightarrow$ $A^{-1} A X = A^{-1} B$
$\Rightarrow$ $\exists!$ решение (1):
\[
  X = A^{-1} B. \tag{3}
\]

? Корректно ли определение? Да, т.\,к. $A^{-1}$ существует для невырожденной
матрицы (определитель $\ne 0$).

\begin{enumerate}[label=\arabic*)]
  \item Действительно, по свойству определителя из (2) $\Rightarrow$
    $\det(A A^{-1}) = \det A \cdot \det A^{-1} = \det E = 1$
    $\Rightarrow$ для существования $A^{-1}$ необходимо
    \[
      \det A \ne 0. \tag{4}
    \]
  \item Допустим $\exists A^{-1}$ $\Rightarrow$ как вычислить:
    \[
      A^{-1} = \frac{1}{\det A} \cdot \widetilde{A}^{\,T}, \tag{5}
    \]
    где $\widetilde{A}$ --- матрица алгебраических дополнений элементов
    матрицы $A$ (присоединённая матрица).
\end{enumerate}

Действительно,
\[
  A^{-1} A = \frac{1}{|A|} \widetilde{A}^{\,T} \cdot A;
\]
строки в левой матрице --- это столбцы алгебраических дополнений для
столбцов $A$.
$\Rightarrow$ В силу замечания в конце предыдущего пункта, при перемножении
$\widetilde{A}^{\,T} A$ по правилу <<строка на столбец>> получим: если
умножаем строку $i$ на столбец $j$, то получим $\det A$ при $j = i$ (то же
разложение по строке $i$) или 0 при $j \ne i$.

В итоге:
\[
  A^{-1} A = \frac{1}{|A|} \begin{pmatrix} |A| & & 0 \\ & \ddots & \\ 0 & & |A| \end{pmatrix} = E.
\]

\subsection{Метод Гаусса для обратной матрицы}

Элементарные преобразования строк:
\begin{enumerate}[noitemsep, label=\arabic*)]
  \item умножение строки на число;
  \item перестановка двух строк;
  \item умножение строки на число и сложение с другой строкой.
\end{enumerate}

Их можно реализовать с помощью левого матричного умножения $P \cdot A$, где
\[
  P = \begin{pmatrix} 1 & & & 0 \\ & \ddots & & \\ & & b & \\ 0 & & & 1 \end{pmatrix}
  \quad (p_{ii} = b)
\]
--- строка $i$ умножится на число $b$.

$P \cdot A$ --- перестановка строк 1, 2, тогда
\[
  P = \begin{pmatrix} 0 & 1 & & 0 \\ 1 & 0 & & \\ & & 1 & \\ 0 & & & \ddots \end{pmatrix}.
\]
? Какой вид $P$ при перестановке произвольных строк.

К строке 1 прибавляется строка 2, умноженная на число $b$. 

? Какой вид $P$
для произвольных строк.

\textbf{Идея:} допустим, в результате элементарных операций над строками
$(A \mid E)$ получилась $(E \mid D)$. Это означает, что существуют матрицы
$P_1, P_2, \dots, P_k$:
\[
  \underbrace{P_k P_{k-1} \dots P_1}_{D} A = E \ \Rightarrow\ D = A^{-1}.
\]

\end{document}
